% Generated draft: requires mathematical and source review.
\documentclass[11pt]{article}
\usepackage[margin=25mm]{geometry}
\usepackage{fontspec}
\setmainfont{Latin Modern Roman}
\usepackage{amsmath,amssymb,mathtools,unicode-math}
\setmathfont{Latin Modern Math}
\usepackage{enumitem,needspace,xurl,hyperref}
\setlength{\parindent}{0pt}
\setlength{\parskip}{5pt}
\setlength{\emergencystretch}{3em}
\newcommand{\N}{\mathbb{N}}
\newcommand{\Z}{\mathbb{Z}}
\newcommand{\Q}{\mathbb{Q}}
\newcommand{\R}{\mathbb{R}}
\newcommand{\C}{\mathbb{C}}
\newcommand{\F}{\mathbb{F}}
\newcommand{\A}{\mathbb{A}}
\newcommand{\PP}{\mathbb{P}}
\newcommand{\E}{\mathbb{E}}
\newcommand{\eps}{\varepsilon}
\newcommand{\dd}{\,\mathrm{d}}
\begin{document}
\section*{UnsolvedMath ID: 800810}
\textbf{Source problem code:} catalogue entry
\subsection{Definitions and mathematical statement}
One conventional interpretation of the notation is the following: if \(M\) is a smooth manifold equipped with a smooth volume form \(\mu\), then
\[
\mathscr{D}_{\mu}(M)=\{\varphi\in\operatorname{Diff}(M):\varphi^{*}\mu=\mu\},
\]
where \(\operatorname{Diff}(M)\) denotes the group of smooth diffeomorphisms of \(M\). This is a conditional explanation of the notation; the supplied record does not specify the manifold, volume structure, regularity, group component, or boundary conditions intended.

For a Riemannian manifold \((X,g)\), local symmetry means that each point \(p\) admits a local isometry \(s_p\) fixing \(p\) with differential \((\dd s_p)_p=-\operatorname{Id}_{T_pX}\). Applying this definition in an infinite-dimensional setting requires a specified manifold structure and Riemannian metric.

\textbf{Source question (Khesin, Misiolek, Modin, 2022, Problem 9).} How ``far'' is \(\mathscr{D}_{\mu}(M)\) from being an infinite-dimensional (locally) symmetric Riemannian space?

The supplied question is qualitative: it specifies neither a Riemannian metric on \(\mathscr{D}_{\mu}(M)\) nor a meaning of ``far.'' Consequently, it does not determine a numerical quantity, a quantified assertion, or a unique criterion to establish. A precise formulation requires these choices and clarification of whether local symmetry, global symmetry, or both are intended.
\subsection{Short English statement}
The question asks how the volume-preserving diffeomorphism group differs from an infinite-dimensional symmetric or locally symmetric Riemannian space. The supplied excerpt leaves the metric, underlying geometric hypotheses, and meaning of ``far'' unspecified.
\subsection{Sources}
\begin{itemize}
\item Local UnsolvedMath registry, permanent ID \(800810\), source code \texttt{catalogue entry}, title: \emph{Geometric Hydrodynamics in Open Problems --- Problem 9}. The registry attributes the problem to Khesin, Misiolek, and Modin, with proposed year 2022, and labels it \texttt{open}; this is a supplied status label, not an independently verified assessment.
\item Khesin, Misiolek, Modin, \emph{Geometric Hydrodynamics in Open Problems} (2022), Problem 9. Supplied source URL: \url{https://arxiv.org/abs/2205.01143}. The supplied provenance records access on 2026-10-16 and extraction by \texttt{arxiv\_latex}; these provenance claims have not been independently verified.
\end{itemize}
\end{document}
